\documentclass[10pt,twocolumn]{article}
\usepackage[T1]{fontenc}
\usepackage[utf8]{inputenc}
\usepackage{lmodern}
\usepackage[margin=1.8cm,top=2.4cm,bottom=2cm,columnsep=0.6cm]{geometry}
\usepackage{fancyhdr}
\usepackage{titlesec}
\usepackage{booktabs}
\usepackage{amsmath,amssymb,amsfonts}
\usepackage{microtype}
\usepackage{xcolor}
\usepackage{mdframed}
\usepackage{parskip}
\usepackage{enumitem}
\usepackage{hyperref}

\definecolor{rulered}{RGB}{180,30,30}
\definecolor{boxgray}{RGB}{245,245,245}
\definecolor{keywordgray}{RGB}{100,100,100}

\pagestyle{fancy}
\fancyhf{}
\fancyhead[L]{\textsc{\small Claude Mag}}
\fancyhead[C]{\small\textit{Machine Learning Research Digest}}
\fancyhead[R]{\small 2026-07-29}
\fancyfoot[C]{\small\thepage}
\renewcommand{\headrulewidth}{0.8pt}

\titleformat{\section}{\normalsize\bfseries\color{rulered}}{}{0pt}{}%
  [\vspace{-4pt}\color{rulered}\rule{\columnwidth}{0.4pt}]
\titlespacing{\section}{0pt}{8pt}{4pt}

\hypersetup{colorlinks=true,urlcolor=rulered,linkcolor=black}

\begin{document}
\setlength{\parindent}{0pt}
\setlength{\parskip}{4pt}

{\small\color{keywordgray}\textbf{Keywords:}
  positional encoding \textbullet{}
  2D-RoPE \textbullet{}
  in-context retrieval \textbullet{}
  transformer architecture}\par\vspace{4pt}

{\LARGE\bfseries\raggedright
  Frontier Language Models Struggle to Copy:
  Text Can Be Better Viewed in 2D}\par\vspace{4pt}

{\large\itshape\raggedright
  A Structural Failure in 1D Positional Encoding Prevents Exact
  String Copying---2D-RoPE Solves It at Training Scales up to 1.4B
  Parameters}\par\vspace{6pt}

{\small\textbf{By} Haodong Wen, Yiran Zhang, Yingfa Chen, Kaifeng Lyu
  (Tsinghua University)}\par\vspace{2pt}

{\small\color{keywordgray}arXiv:2607.16072 \textbullet{} July 15, 2026}
\par\vspace{2pt}\noindent{\color{rulered}\rule{\columnwidth}{0.6pt}}\vspace{6pt}

Frontier language models solve complex multi-step reasoning yet
consistently fail at a task any human child can perform: copying a
string exactly. A new paper from Tsinghua University traces this failure
to a structural flaw in standard 1D rotary positional encodings and
proposes \textbf{2D-RoPE}, which organizes tokens into a two-dimensional
grid so that copying reduces to zero-offset column retrieval.
Shallow Transformers with 2D-RoPE achieve perfect copying on sequences
hundreds of times longer than their training range.

\begin{mdframed}[backgroundcolor=boxgray,linecolor=rulered,linewidth=1pt,
    innerleftmargin=6pt,innerrightmargin=6pt,
    innertopmargin=6pt,innerbottommargin=6pt]
\textbf{\small AT A GLANCE}\par\vspace{4pt}
\begin{description}[leftmargin=0pt,labelsep=4pt,itemsep=2pt,parsep=0pt,
    font=\small\bfseries]
  \item[Problem:] Standard 1D positional encodings force copying to
    rely on local-context shortcuts whose relative offsets change with
    sequence length, causing systematic failure on long strings.
  \item[Why it matters:] Verbatim extraction from context---for
    quotes, code formatting, structured outputs---is foundational; a
    structural fix is more reliable than prompt engineering.
  \item[Method:] 2D-RoPE assigns each token a row ID and column ID;
    copying at fixed column offset becomes a zero-relative-offset
    attention task.
  \item[Result:] 6-layer Transformers with 2D-RoPE reach 99.7\%
    exact copy accuracy at lengths 8$\times$ beyond training range;
    standard RoPE achieves 12\%.
  \item[Evidence:] Advantage holds in 1.4B pretraining experiments
    and on downstream verbatim extraction benchmarks.
\end{description}
\end{mdframed}

\section{The Big Picture}

Large language models have surpassed humans on olympiad mathematics,
bar exams, and competitive programming---but they cannot reliably copy
a 1024-token string they just read. This asymmetry is not a quirk of
training data; it reflects a structural mismatch between how attention
mechanisms encode position and what the copy task requires.

The paper demonstrates the failure empirically on GPT-4o, Claude 3.7,
and Gemini 2.0: exact-copy accuracy falls below 50\% on 512-token
strings and degrades further on repetitive inputs. The authors then
derive a root cause from the mathematics of RoPE, and propose a fix
that achieves perfect copying at lengths orders of magnitude beyond
training.

\section{Why Existing Approaches Fall Short}

In Transformer attention with RoPE, the attention score between query
position $p$ and key position $q$ depends on the relative offset $p - q$.
Copying token $x_k$ to output position $y_k$ requires attending to
exactly input position $k$. But as the output grows, the relative offset
$k - y_k$ changes---so the model must infer $k$ dynamically rather
than attend to a fixed offset.

In practice, models learn to copy by matching local context: if a
substring appears only once, the model retrieves it via pattern matching
rather than positional address. This works for unique strings but breaks
when strings repeat (the context match is ambiguous) or are very long
(the attention score for the correct position degrades relative to
spurious matches). Prompt-engineering workarounds (``copy exactly'')
improve but do not solve the problem because the architectural inductive
bias remains.

\section{Core Mechanism}

2D-RoPE organizes a sequence of $N$ tokens into a grid with $R$ rows and
$C = N/R$ columns. Token $i$ receives:
\begin{itemize}[itemsep=1pt,parsep=0pt]
  \item Row ID: $r_i = \lfloor i/C \rfloor$
  \item Column ID: $c_i = i \bmod C$
\end{itemize}

Rotary encoding is applied along each dimension independently. Under this
view, copying is: for output position $(r, c)$, retrieve the input token
at the same column $c$ in row $r$. The column offset between any output
token and its corresponding input token is always zero---a fixed,
length-independent offset that the model can reliably learn.

The key invariant: \textit{the relative column offset for copying is 0,
regardless of sequence length}. Standard 1D RoPE has no equivalent fixed
offset.

\section{Technical Formulation}

Standard RoPE applies rotation $e^{ip\theta}$ at position $p$ with
frequency $\theta$, giving attention score:
\[
  A_{pq} = \mathrm{Re}\!\left[ q_p^* k_q \cdot e^{i(p-q)\theta} \right]
\]

2D-RoPE splits the embedding dimension: half the heads attend to
row position, half to column position:
\[
  A_{ij} = \mathrm{Re}\!\left[ q_i^* k_j \cdot e^{i\left[(r_i - r_j)\theta_R
    + (c_i - c_j)\theta_C\right]} \right]
\]

During a copy task, $r_i - r_j = 0$ (same row, input vs.\ output)
and $c_i - c_j = 0$ (same column), so:
\[
  A_{ij}^\text{copy} = \mathrm{Re}\!\left[ q_i^* k_j \right]
\]

The attention score reduces to the inner product of query and key
without any positional modulation---a fixed, position-independent signal
that is identical at any sequence length.

\section{Learning or Inference Procedure}

Training proceeds identically to standard causal language modeling,
replacing the 1D position index with the 2D $(r, c)$ pair. No changes
to the loss function, optimizer, or data pipeline are required. The grid
dimensions $R$ and $C$ are treated as hyperparameters; the paper uses
$C = 64$ for all experiments (so rows contain 64 tokens).

At inference, sequences longer than training length are handled by
extending the row dimension: a 4096-token sequence with $C = 64$
requires $R = 64$ rows, while a 512-token training sequence requires
only $R = 8$ rows. Column IDs remain bounded by $C$, so the column
encoding generalizes trivially to longer sequences.

\section{What the Guarantee Says}

For a sequence of length $N$ with column width $C$, 2D-RoPE guarantees
that the copy attention signal is \emph{position-independent} at offset
$(0, 0)$ in row-column space. This is a structural, not statistical,
guarantee: no amount of sequence length increase can alter the positional
signal the model must attend to in order to copy correctly.

The guarantee does not cover sequences longer than $R_\text{max} \cdot C$
(where $R_\text{max}$ is the training row count), but row generalization
is analogous to standard length generalization and can be addressed with
existing techniques (ALiBi, YaRN).

\section{Experimental Findings}

\textbf{Synthetic copy task} (exact accuracy at length $L$, trained on $L=512$):

\begin{center}\small
\begin{tabular}{lcc}
\toprule
Model & $L=2048$ & $L=4096$ \\
\midrule
Standard RoPE (6L) & 34.1 & 12.3 \\
2D-RoPE (6L) & 98.4 & 99.7 \\
Standard RoPE (12L) & 41.8 & 18.7 \\
2D-RoPE (12L) & 99.1 & 99.8 \\
\bottomrule
\end{tabular}
\end{center}

\textbf{Repetitive strings} (period 4, length 512, in-distribution):
Standard RoPE: 7.3\%. 2D-RoPE: 98.1\%.

\textbf{Pretraining at 1.4B parameters} (100B tokens, same perplexity
as baseline): 2D-RoPE models improve on verbatim extraction benchmarks
by 14.2 points while matching or exceeding baseline on MMLU, HellaSwag,
and ARC-Challenge.

\section{Ablations and Interpretation}

\textbf{Grid aspect ratio:} Varying $C \in \{16, 32, 64, 128\}$ shows
$C = 64$ is optimal; very narrow grids ($C = 16$) reduce copy benefit by
requiring more rows and increasing within-row attention competition.

\textbf{Attention visualization:} On copy tasks, 2D-RoPE heads in the
column-position half show clean offset-0 attention spikes at the source
token; standard RoPE models show diffuse attention over a $\pm 5$ token
window, confirming the local-context shortcut.

\textbf{Frontier model analysis:} Attention rollout on GPT-4o and
Claude 3.7 during copy failures reveals the same diffuse-attention
signature, confirming the failure mode is not specific to small models.

\textbf{Column vs.\ row only:} Using only column encoding (dropping row)
degrades general language modeling perplexity by 0.4 bits while
maintaining near-perfect copying---confirming that row and column
encodings serve distinct roles.

\section*{Reference}

Haodong Wen, Yiran Zhang, Yingfa Chen, Kaifeng Lyu.
\textbf{Frontier Language Models Struggle to Copy: Text Can Be Better Viewed in 2D}.
arXiv:2607.16072, July 15, 2026.
\url{https://arxiv.org/abs/2607.16072}

\end{document}
